
\subsection{Special case of AMSS model}\label{sec:specialcaseAMSS}%
That the  Ramsey allocation for the AMSS model differs from the Ramsey allocation of the Lucas-Stokey model is a symptom that the measurability
constraints \Ep{eqn:AMSSapp6} bind.
Following Bhandari, Evans, Golosov, and Sargent (2017) (henceforth BEGS), we now consider a special case of the AMSS model in which these constraints don't bind.
Here  the AMSS Ramsey planner chooses not to issue state-contingent debt, though he is free to do so.  The environment is one in which
fluctuations in the risk-free interest rate provide all  the insurance that the Ramsey planner wants.

Following BEGS, we set $S=2$ and assume that
the state $s_t$ is i.i.d., so that the transition matrix $\Pi(s'|s) = \Pi(s')$ for $s=1,2$.
Following BEGS, it is useful to consider  the following special case of the implementability constraints  \Ep{eqn:AMSSapp6} evaluated at the constant
value  of the state variable  $x_-= x(s) = \check x$:
$$ {\frac{u_c(s) \check x}{\beta \sum_{\tilde s} \Pi(\tilde s) u_c(\tilde s) }} = u_c(s) (n(s) - g(s) ) - u_l(s) n(s) + \check x, \ s=1,2. \EQN eqn_AMSS_app_100 $$
We   guess and verify that the scaled Lagrange multiplier $\mu(s)=\mu$ is a constant independent of $s$. At a fixed $x$,  because $V_x(x, s)$ must be independent of $s_-$, equation \Ep{eqn:AMSSapp9} becomes
$ V_x(\check x) = \sum_{s} \left( \Pi(s) {\frac{u_c(s)}{\sum_{\tilde s} \Pi(\tilde s) u_c(\tilde s)}} \right) V_x(\check x) = V_x(\check x) ,$%QN eqn:AMSSapp109 $$
which confirms that equation \Ep{eqn:AMSSapp7} becomes  $\mu = \beta V_x(\check x)$ for a  constant  $\mu$.
For the continuation Ramsey planner facing implementability constraints \Ep{eqn_AMSS_app_100}, the first-order conditions with respect to $n(s)$ become
$$ \EQNalign{ & u_c(s) - u_l(s)  \mu \Bigl\{ {\frac{\check x}{\beta \sum_{\tilde s} \Pi(\tilde s) u_c(\tilde s)}} (u_{cc}(s) - u_{cl}(s) ) - u_c(s) \cr
       &  \ \ - n(s) (u_{cc}(s) - u_{cl}(s)) - u_{lc}(s) n(s) + u_l(s) \Bigr\}  = 0.\EQN eqn:AMSSapp101 }  $$
Equations  \Ep{eqn_AMSS_app_100} and \Ep{eqn:AMSSapp101}  are four equations in the four  unknowns $n(s), s=1,2$, $\check x$, and $\mu$. Under some regularity conditions
on  the period utility function  $u(c,l)$, BEGS show that these equations have a unique solution featuring a negative value of $\check x$.
Consumption $c(s)$ and the flat-rate tax on labor $\tau(s)$ can then be constructed
as history-independent functions of $s$.
\auth{Evans, David G.}%
\auth{Bhandari, Anmol}%
\auth{Golosov, Mikhail}%
\auth{Sargent, Thomas J.}%
In this AMSS economy, $\check x = x(s) = u_c(s) {\frac{b_{t+1}(s)}{R_t(s)}}$.  The risk-free interest rate, the tax rate, and the marginal
utility of consumption fluctuate with $s$, but $x$ does not and neither does  $\mu(s)$. The labor tax rate and
the allocation depend only on the current value of $s$.

For this special AMSS economy to be in a steady state from time $0$ onward, it is necessary that  initial debt $b_0$ satisfy the
time $0$ implementability
constraint \Ep{eqn:AMMSSapp101} at the value $\check x$ and the realized value of $s_0$.  We can solve for this level of $b_0$ by
plugging the $n(s_0)$ and $\check x$
that solve our four equation system  \Ep{eqn_AMSS_app_100}-\Ep{eqn:AMSSapp101}  into
$ u_{c,0} b_0 = u_{c,0} (n_0-g_0) - u_{l,0} n_0 + \check x $ and solving for $b_0$.
This  $b_0$ assures that a steady state $\check x$ prevails from time $0$ on.


\subsubsection{Relationship to a Lucas-Stokey economy}
The constant value $\mu$ of the Lagrange multipliers $\mu(s)$ across states $s$  in the Ramsey plan for our special AMSS economy is a tell tale sign that
the measurability restrictions imposed on the Ramsey allocation by the requirement that government debt must be risk free are slack.
When they bind, those
measurability restrictions  cause the AMSS tax policy and allocation to be history dependent -- that's what activates flucations in the risk-adjusted martingale.
Because  those measurability conditions are  slack in  this special AMSS economy, the same Ramsey allocation is attained in a corresponding complete markets
 a Lucas-Stokey economy that starts
from a particular initial  government debt.  The setting of  $b_0$ for the corresponding Lucas-Stokey
implementability constraint \Ep{Bellman2cons} solves
$ u_{c,0} b_0 = u_{c,0}(n_0 - g_0) - u_{l,0} + \beta \check x$.
In this Lucas-Stokey economy, although the Ramsey planner is free to issue state-contingent debt, it chooses not to and instead  issues  only risk-free debt.
It achieves the
 risk-sharing with the private sector that it wants by altering the amounts of one-period risk-free debt that it issues at each current state,
  while taking advantage of  the fluctuations in the  equilibrium interest
rate that its tax policy induces.

\subsubsection{Convergence to the special case}
In an  i.i.d., $S=2$  AMSS economy in which  the initial $b_0$ does not equal the special value described in the previous subsection, $x$ fluctuates and is history-dependent.
%namely, a value at the implied stationary distribution of risk free government debt
%affiliated with the fixed  $\check x$ determined above,
The Lagrange multiplier $\mu_s(s^t)$ is a non trivial risk-adjusted martingale and the allocation and distorting tax rate
are both  history dependent.  However, BEGS describe conditions  under which such an i.i.d., $S=2$ AMSS economy in which
the representative agent dislikes consumption risk  converges to a
Lucas-Stokey economy in the sense that $x_t \rightarrow \check x$ as $t \rightarrow \infty$.
The  following subsection  displays a numerical example that exhibits convergence.



\subsection{Computed examples}%\label{sec:AMSSexamples}}

Figures \Fg{AMSS1}, \Fg{AMSS2}, and \Fg{AMSS3} report features of a Ramsey plan for an AMSS economy with
one-period utility function
$$ u(c,l) = {\frac{c^{1-\sigma}}{1-\sigma}} - {\frac{l^{1+\gamma}}{1+\gamma}} . $$
We set $\beta = .96, \gamma = 2$, and $\sigma$ equal to either $0$ (this is the
quasi-linear case of AMSS) or $1.5$ (here the representative consumer is averse to
consumption fluctuations).  We set
$S=2$ with $g(s_1) = .05$ and $g(s_2) = .15$ and assume that $s_t$ is an iid process with
the probability of both states equalling $.5$.  We computed the objects shown in figures \Fg{AMSS1}--\Fg{AMSS3}
by numerically solving  Bellman equations \Ep{eqn:AMSSapp5} and \Ep{eqn:AMSSapp6}.


\midfigure{AMSS1}
\centerline{\epsfxsize=3true in\epsffile{AMSS1.eps}}
\caption{Two AMSS economies. Dashed line: $\sigma=0$; solid line: $\sigma = 1.5$. Outcomes converge to those associated with  first-best allocation for the
quasi-linear ($\sigma =0$) economy, while they converge to an allocation associated with a complete markets economy for the economy with risk-aversion ($\sigma=1.5$).}
\infiglist{AMSS1}
\endfigure


\midfigure{AMSS2}
\centerline{\epsfxsize=3true in\epsffile{AMSS2.eps}}
\caption{Decision rule for $x_{t+1}- x_t$ as a function of $x_t$ and $s_t$ for AMSS $\sigma=0$ economy.
Dashed line: $g_t$ is high; solid line, $g_t$  is low.}
\infiglist{AMSS2}
\endfigure



\midfigure{AMSS3}
\centerline{\epsfxsize=3true in\epsffile{AMSS3.eps}}
\caption{Decision rule for $x_{t+1}- x_t$ as a function of $x_t$ and $s_t$ for AMSS $\sigma=1.5$ economy.
Dashed line: $g_t$  is high; solid line, $g_t$  is low. }
\infiglist{AMSS3}
\endfigure

%
%
Figure \Fg{AMSS1} reports outcomes from
long simulations of the two AMSS economies, the quasi-linear ($\sigma=0)$ economy shown in the dashed line and the
 $\sigma=1.5$ economy in the solid line.  As expected from the AMSS results summarized in section \use{sec:AMSS} of chapter \use{optax}, outcomes in the  quasi-linear economy
converge to those
associated with a first-best allocation. The government accumulates sufficient risk-free claims on the private sector  to finance all subsequent government expenditures
with interest earned from those claims. After accumulating enough assets,
the government sets the tax rate to zero and rebates excess earnings
from assets in the form of  lump sum transfers to the representative agent.

Asymptotic outcomes  for the $\sigma = 1.5$ economy display the striking outcome that they too converge, but not to a first-best allocation.  Instead, they converge to the
fixed point in $x$  associated with a Lucas-Stokey complete markets allocation.

Figure \Fg{AMSS2} displays the Ramsey plan decision rule for $x_{t+1}$  for the $\sigma=0$ quasi-linear AMSS economy,  while \Fg{AMSS3} displays
 the decision rule for the $\sigma =1.5$ economy.  The figures display $x_{t+1} - x_t$ as a function of $x_t$ and $g_t$.   Rest points  for $x$ occur where the branch of
the decision rule affiliated with  high $g$ intersects the branch affiliated with low $g$ at  $x_{t+1} - x_t$ equal to zero.
For the quasi-linear economy, the branches intersect at a value of  $x$ and an implied
level of government debt that is sufficient to finance all subsequent government expenditures. Further, when
$x_t$ exceeds this level,  the decision rules on average point downward, imparting a negative drift to $x_t$ until it reaches the level associated with enough
government assets capable of sustaining a first-best allocation.  Notice that after $x_t$ reaches this level,  the tax rate equals zero.


For the figure \Fg{AMSS3} AMSS economy with $\sigma=1.5$, the branches of the decision rule for $x_{t+1} -x_t$ intersect at coordinate  value  zero at an ordinate value of $x$ that,
while negative, exceeds the
level for the quasi-linear  economy.  This $x$ is affiliated with a Lucas-Stokey economy.  After $x_t$ has converged to this level,
the tax rate is positive and constant. That it does not fluctuate
with $g$ is partly a consequence of the specification of the one-period utility function.
